\chapter{Background}\label{ch:background}

In this chapter we address the building blocks for the development of our guidance method.
We start with a short introduction to flow matching, the generative framework used by our model of choice.
We then proceed by summarizing the idea of guidance and expose a few selected methods.
We subsequently outline the relevant implementation details of ArchesWeather, and conclude with a short literature review on the use of guidance in the context of generative weather modeling. 

\section{Flow matching}

Generative modeling frames the task of generating objects $x \in \mathbb{R}^d$, for instance images or text, as sampling from a data distribution $p^1$.
In particular, flow matching \parencite{lipman2023flow} models the data distribution $p_1$ as the endpoint of a
continuous transport of a tractable prior $p_0$, via the probability path $(p_t)_{t\in[0,1]}$.
This path is generated by a time-dependent velocity field $u_t$, defining an ordinary differential equation (ODE), whose solution is the flow $\phi_t$:
$$
    \frac{\mathrm{d}}{\mathrm{d}t}\,\phi_t(x) = u_t\big(\phi_t(x)\big), \qquad \phi_0(x) = x.
$$
In other words, the push-forward of $\phi_t$ carries $p_0$ to $p_t$.
% ; equivalently, $u_t$ and $p_t$ obey the continuity equation $\partial_t p_t + \nabla\!\cdot(p_t\,u_t) = 0$.
% Sampling then reduces to drawing $x_0 \sim p_0$ and integrating the flow to $t=1$. Generally the ODE has no closed form solution and is thus simulated numerically, for example via Euler's method.

A common choice for $(p_t)_{t\in[0,1]}$ is a Gaussian conditional path, 
which uses two noise schedulers $\alpha_t, \beta_t \in [0,1]$ to interpolate a data sample $x_1$ with noise:
$$
    p_t(\cdot \mid x_1) = \mathcal{N}\!\big(\alpha_t\, x_1,\ \beta_t^{2}\, I\big)
$$
The schedulers are chosen such that $\alpha_t$ rises from $0$ to $1$ while $\beta_t$ falls from $1$ to $0$ such that 
\begin{align*}
    p_0(\cdot|x) &= \mathcal{N}(\mathbf{0}, \mathbf{1}) \\
    p_1(\cdot|x) &= \delta_x
\end{align*}

Samples from the marginal density $p_t$ can then be easily obtained by drawing a data sample $x_1 \sim p_1$ and noise $x_0 \sim \mathcal{N}(0,I)$, and interpolate them with
$$
    x_t = \alpha_t\, x_1 + \beta_t\, x_0 \ \sim\ p_t .
$$
To generate samples that follow this probability path, in flow matching we train a neural network $u_t^\theta$ set to generate $p_t$. 
This is obtained by noting that, although the marginal field is
not immediately available to us, it decomposes as an expectation of conditional fields
$$u_t(x) = \mathbb{E}_{x_1 \sim p_{1\mid t}(\cdot\,\mid x)}\!\big[\,u_t(x \mid x_1)\,\big]$$
each $u_t(\cdot \mid x_1)$ generating a simple path from $p_0$ to a fixed data sample $x_1$. Flow matching
exploits the fact that the tractable conditional objective
\begin{equation}
    \mathcal{L}(\theta) = \mathbb{E}_{t,\; x_1 \sim p_1,\; x_t \sim p_t(\cdot\,\mid x_1)}
    \big\lVert u_\theta(x_t, t) - u_t(x_t \mid x_1) \big\rVert^2
\end{equation}
shares its gradient with the intractable marginal one \parencite{lipman2023flow}. Under the linear
(optimal-transport) path $x_t = (1-t)\,x_0 + t\,x_1$ with $x_0 \sim p_0$ ($\alpha_t=t$, $\beta_t=1-t$), the conditional field is the
constant displacement $u_t(x_t \mid x_1) = x_1 - x_0$, such that training amounts to a least-squares
regression onto $x_1 - x_0$.

---somewhere say that since we are drawing from noise we can generate independent samples.

\section{Guidance methods}
Guidance methods were originally develped in the context of diffusion models trained to generate images.
The first among these ideas is classifier guidance.
In the following we make a short derivation of the method, translate it to the flow matching framework, and expose later perspectives that were developed after the first ideas. 

\subsection{Classifier guidance}
In most cases, we are not interested in sampling an arbitrary object from $p^1$, but we rather to condition the generation on a certain property $y$.
Classifier guidance builds on the idea that guided generation can be framed as sampling from the posterior distribution $p^1(x|y)$. 
To do so, we have to specify the conditional vector field $u_t^\theta(z_t|y)$.
Using Bayes rule we can derive the following:
\begin{align*}
p(z_t \mid y) &\propto p(y \mid z_t)\, p(z_t), \\
\nabla_{z_t}\log p(z_t \mid y)
&= \nabla_{z_t}\log p(z_t)
+ \nabla_{z_t}\log p(y \mid z_t).
\end{align*}
Using the conversion formula between diffusion and flow matching, and inserting the previous result, we can express the conditional vector field as
\begin{align*}
    u_t^\theta(z_t|y) &= a_t \nabla_{z_t} \log p_t(z_t|y) + b_t z_t \\
    &= u_t^\theta(z_t) + a_t \nabla_{z_t} \log p_t(y|z_t)
\end{align*}
The guided vector field can thus be expressed as a sum between the unguided vector field and the gradient with respect to the noisy state of the likelihood $p_t(y|z_t)$.

Originally, to obtain guided samples researchers used the gradients of an image classifier, hence the name of the method.

\subsection{Guidance strength-schedule}

While the latter mathematical framework offers an elegant way to express guidance as sampling from the posterior distribution $p^1(x|y)$,
researchers realised that they could obtain better results by scaling up the contribution of the guidance term.

To unify the notation we can express this idea as follow: replace $a_t$ with a general guidance strength-schedule $\lambda_t$,
which we can decompose 
\begin{equation}
    \lambda_t := w \cdot \alpha_t \cdot c_t
\end{equation}
with:
\begin{itemize}[noitemsep]
    \item $w$: guidance strength (time independent)
    \item $\alpha_t$: guidance profile
    \item $c_t$: a general step normalization constant (specific to different guidance algorithms)
\end{itemize}
Under this decomposition we would retrieve the approximate posterior sampling by setting $w=1$, $\{\alpha_t\}_{t=1}^T=\{a_t\}_{t=1}^T$, and $c_t=1$.
Note that for $\lambda_t \neq a_t$ for $t=1,\dots, T$, the guided results stops being a sample from the posterior distribution.
In this sense we trade mathematical reasoning for control.

Different approaches to scale guidance have been to:
\begin{itemize}[noitemsep]
    \item scale up $w$
    \item change the default profile schedule to an arbitrary $\alpha_t$
    \item normalize the step to satisfy a certain constraint, such as setting $\{\lambda_t\}_{t=1}^T$ s.t. the guidance vector represents a fix fraction of the unguided vector field.
\end{itemize}

The question for which guidance schedule is most optimal remains unanswered.

Applying this to our context:
\begin{align*}
    \tilde{u}(z_t) &= u(z_t) - a_t \nabla_{z_t} \log p(y|z_t) \\
    &= u(z_t) - \lambda_t \nabla_{z^t}\mathcal{L}(\hat{x}_t)
\end{align*}

In \ref{...} develop variants of the base algorithm we will present next, testing empirically appropriate choices for $\{\lambda\}_{t=1}^T$ from different perspectives.


---The quest for the "best schedule" is still something that hasn't been solved in the guidance literature, although many heuristics have been proposed.


\subsection{Backpropagating through the ODE}

---Since this map from $x_0$ to $x_1$ is a single deterministic and differentiable trajectory, its output can be steered by differentiating a
target loss through the flow---the mechanism on which our guidance builds (\Cref{ch:guide}).

\subsubsection{FlowGrad}

\subsubsection{D-Flow}

\subsubsection{Control paper}

FlowGrad \parencite{liu2023flowgrad} backpropagates such a gradient through the sampling trajectory; D-Flow
\parencite{benhamu2024dflow} instead optimises the source $x_0$, which for Gaussian paths projects
the update onto the data manifold; and \textcite{wang2025trainingfree} frame the steering as optimal
control. Our method adopts this gradient-based view but, rather than tuning a guidance strength,
drives an explicit target gap to zero (\Cref{ch:guide}).


\section{ArchesWeather}

Since the start of the machine learning weather literature a few models based on generative architectures have been proposed.
Among the most prominent ones we can find GenCast.

GenCast ... operational ...

To prototype our method however the computational and memory requirements for GenCast are quite elevate.
For this reason we pick as model of choice ArchesWeather.

We now spend some time presenting ArchesWeather in detail.
This will enable us to introduce the relevant pieces we will be working on.

In our setting the model is a
\emph{residual} predictor: the generated $x_1$ is an increment added to a base state (a previous
state or a deterministic forecast), so each ensemble member is that base plus one integrated
trajectory (\Cref{ch:guide}).

\subsection{Dataset}

\subsection{Architecture}

\subsection{Training}

\subsection{Evaluation}

\section{Guidance in weather modeling}

The last piece of the puzzle before presenting our approach is to expose what people have alredy tried out, enabling us to expose the relevant gaps that we want to fill.

Digging into paper archives, I was able to find only three notable papers that were approaching the same general task as us.

The first one is ...

The most similar in nature, and the one from which we will blatently borrow some aspects of the evaluation, is also the one that approaches guidance and its components the most thorrouwly.

We have thus presented all relevant background pieces for the exposition of our contributions.
We will spend the next three chapters introducing each of these key aspects. 